\section{Extreme stabilization times}

For $M\ge2$, let
\begin{equation}\label{eq:bM}
  b_M=8\log M-4\log\log M-4\log\pi.
\end{equation}
All logarithms are natural.

\begin{theorem}[Gumbel law and Poisson terminal defects]\label{thm:gumbel}
Let $T_1,\dots,T_M$ be the stabilization marks on any $M$ prescribed
coordinates and let $T_M^{\max}=\max_{1\le i\le M}T_i$.  Then, for every
$x\in\R$,
\begin{equation}\label{eq:gumbel-limit}
  \P\!\left(\frac{T_M^{\max}-b_M}{8}\le x\right)
  \longrightarrow \exp(-e^{-x}).
\end{equation}
At the same scale, the remaining terminal-defect count satisfies
\begin{equation}\label{eq:poisson-defect-limit}
  N_M(b_M+8x)
  \ \xrightarrow{\mathrm d}\
  \operatorname{Pois}(e^{-x}),
\end{equation}
where $N_M(t)=\sum_{i=1}^M\1\{T_i>t\}$.

The convergence in \eqref{eq:gumbel-limit} holds in every fixed moment.  In
particular,
\begin{align}
  \E T_M^{\max}
  &=b_M+8\gamma+o(1),\label{eq:mean-max}\\
  \Var(T_M^{\max})
  &=\frac{32\pi^2}{3}+o(1),
  \label{eq:variance-max}
\end{align}
where $\gamma$ is Euler's constant.
\end{theorem}

\begin{proof}
The normal-tail expansion gives
\begin{equation}\label{eq:q-asymptotic}
  q(t)
  =\sqrt{\frac{8}{\pi t}}e^{-t/8}
   \left(1-\frac4t+O(t^{-2})\right)
  \qquad(t\to\infty).
\end{equation}
For fixed $x$, put $t_M(x)=b_M+8x$.  From \eqref{eq:bM},
\begin{equation}\label{eq:centering-identity}
  M e^{-b_M/8}=\sqrt{\pi\log M}.
\end{equation}
Combining \eqref{eq:q-asymptotic} and \eqref{eq:centering-identity} yields
\begin{equation}\label{eq:tail-intensity-limit}
  Mq(t_M(x))\longrightarrow e^{-x}.
\end{equation}
Since $q(t_M(x))\to0$, the exact formula
\[
  \P(T_M^{\max}\le t_M(x))=(1-q(t_M(x)))^M
\]
gives \eqref{eq:gumbel-limit}.  The exact binomial law
\[
  N_M(t_M(x))\sim\operatorname{Bin}(M,q(t_M(x)))
\]
and \eqref{eq:tail-intensity-limit} give
\eqref{eq:poisson-defect-limit}.

For completeness, we verify the uniform integrability needed for moment
convergence.  Mills-ratio bounds imply constants $0<c<C<\infty$ such that
\begin{equation}\label{eq:two-sided-q-bound}
  c\frac{e^{-t/8}}{\sqrt{1+t}}
  \le q(t)\le
  C\frac{e^{-t/8}}{\sqrt{1+t}},
  \qquad t\ge0.
\end{equation}
Write $V_M=(T_M^{\max}-b_M)/8$.  For $x\ge0$ and all sufficiently large $M$,
\begin{equation}\label{eq:right-tail-UI}
  \P(V_M>x)
  \le Mq(b_M+8x)
  \le C_1e^{-x}.
\end{equation}
For $0\le x\le b_M/16$, the lower bound in
\eqref{eq:two-sided-q-bound} gives
\begin{equation}\label{eq:left-tail-UI-one}
  \P(V_M<-x)
  \le \exp\{-Mq(b_M-8x)\}
  \le \exp\{-c_1e^x\}.
\end{equation}
For $b_M/16\le x\le b_M/8$, monotonicity gives
\begin{equation}\label{eq:left-tail-UI-two}
  \P(V_M<-x)
  \le \exp\{-Mq(b_M/2)\}
  \le \exp\{-c_2M^{1/2}/\sqrt{\log M}\}.
\end{equation}
For $x>b_M/8$, the probability is zero because stabilization times are
nonnegative.  The bounds \eqref{eq:right-tail-UI}--\eqref{eq:left-tail-UI-two}
imply uniform integrability of every fixed power of $|V_M|$.  Hence the moments
converge to those of a standard Gumbel random variable, whose mean and variance
are $\gamma$ and $\pi^2/6$, respectively.  This proves
\eqref{eq:mean-max}--\eqref{eq:variance-max}.
\end{proof}

\begin{remark}[Interpretation of the centering]\label{rem:centering}
The term $8\log M$ comes from the exponential factor $e^{-t/8}$ in the gamma
tail.  The correction $-4\log\log M$ comes from its $t^{-1/2}$ prefactor, and
$-4\log\pi$ fixes the limiting Poisson intensity to be $e^{-x}$.
\end{remark}
